%% 4_experiments
\section{What Limits Learned Reconstruction}
\label{sec:experiments}

\subsection{One protocol for every number}

Every number in this section comes from one protocol, under which all methods see the
same scenes, the same trajectories and the same held-out cells, and are trained to
the same budget. We evaluate on RadioMapSeer, a ray-traced corpus of urban scenes with dense
ground truth, using $40$ held-out maps and $8$ transmitters each for $320$ test
scenes. Every method sees the same trajectories, the same measurement noise and the
same held-out cells. We report root mean square error over free-space cells that
were never measured, with bootstrap $95\%$ confidence intervals over $2000$
resamples and paired Wilcoxon signed-rank tests over the $320$ scenes. Deep
baselines are trained to the same budget of $150$ epochs on the same split.

Holding that budget fixed forces a decision about what a baseline is, and we resolve
it by separating faithful reimplementations from matched-budget representatives of an
architecture family, a distinction \cref{tab:main} marks throughout. RadioUNet, PMNet, RadioDiff and RMDM are reimplementations of a
specific published architecture, following the paper in structure and differing only
where our input setting requires it. The entries marked with a dagger are not
reimplementations of one paper. They are matched-budget representatives of an
architecture family, built to the same conditioning interface as everything else so
that the family is the only variable, and the citation beside each names the work
that family is drawn from rather than a method we reproduced. We prefer this for the
question at hand. Comparing author implementations across families would confound
the backbone with six different training recipes, whereas the claim we are testing is
precisely whether the family matters once everything else is held fixed.

Protocol discipline matters more in this literature than in most, because the
conventions in circulation disagree by enough to reverse a ranking. The same
checkpoint scores between $8.6$ and $9.5$~dB under the four evaluation conventions
we found in use, a spread of nearly a decibel on a scale where the leading methods
are separated by less, so a table mixing conventions would carry no information.

\begin{table}[t]
\centering\scriptsize
\setlength{\tabcolsep}{3pt}
\begin{tabular}{@{}llrrrr@{}}
\toprule
& Method & RMSE$\downarrow$ & SSIM$\uparrow$ & LoS$\downarrow$ & NLoS$\downarrow$ \\
\midrule
\multicolumn{6}{@{}l}{\textit{Ours}} \\
& WNet $+$ occlusion & \textbf{8.64}\,\tiny[8.41,8.85] & 0.841 & 3.85 & 9.78 \\
& U-Net $+$ occlusion & 9.47\,\tiny[9.20,9.72] & 0.826 & 3.95 & 10.73 \\
\addlinespace[1pt]
\multicolumn{6}{@{}l}{\textit{Learned baselines}} \\
& RadioUNet~\cite{levie2021radiounet} & 9.93\,\tiny[9.68,10.18] & 0.816 & 5.28 & 11.12 \\
& WNet (no occlusion) & 10.01\,\tiny[9.72,10.28] & 0.825 & 5.18 & 11.26 \\
& RadioMamba & 10.92\,\tiny[10.65,11.19] & 0.795 & 7.00 & 12.10 \\
& RadioTransf. & 10.97\,\tiny[10.66,11.28] & 0.791 & 6.53 & 12.19 \\
& URAM & 11.37\,\tiny[11.11,11.61] & 0.801 & 7.72 & 12.50 \\
& RadioDiff & 11.37\,\tiny[11.03,11.69] & 0.803 & 5.99 & 12.74 \\
& RadioGAN & 11.42\,\tiny[11.15,11.67] & 0.788 & 7.09 & 12.63 \\
& U-Net (no occlusion) & 11.51\,\tiny[11.17,11.85] & 0.797 & 6.46 & 12.84 \\
& PMNet~\cite{lee2024pmnet} & 14.12\,\tiny[13.71,14.50] & 0.714 & 9.98 & 15.52 \\
& SparseUNet & 21.05\,\tiny[20.67,21.42] & 0.656 & 26.89 & 20.17 \\
& RMDM & 35.87\,\tiny[34.41,37.25] & 0.715 & 11.19 & 39.67 \\
\addlinespace[1pt]
\multicolumn{6}{@{}l}{\textit{Input ablations}} \\
& Tx-free occlusion & 9.88\,\tiny[9.59,10.17] & 0.825 & 4.96 & 11.16 \\
& no Tx heatmap & 16.25\,\tiny[15.80,16.71] & 0.739 & 19.44 & 16.32 \\
& no buildings & 18.82\,\tiny[18.46,19.19] & 0.674 & 20.84 & 18.82 \\
\addlinespace[1pt]
\multicolumn{6}{@{}l}{\textit{Classical}} \\
& RBF & 25.64\,\tiny[24.85,26.44] & 0.720 & 27.19 & 26.42 \\
& Ord. Kriging & 26.99\,\tiny[26.40,27.59] & 0.693 & 34.25 & 26.68 \\
& IDW ($p{=}1$) & 27.85\,\tiny[27.07,28.60] & 0.699 & 27.80 & 28.92 \\
& GP & 28.09\,\tiny[27.47,28.71] & 0.674 & 35.13 & 27.75 \\
\addlinespace[1pt]
\bottomrule
\end{tabular}
\caption{\textbf{Free-space unobserved RMSE (dB) on $320$ held-out scenes},
with bootstrap $95\%$ intervals, structural similarity, and the error split by
line-of-sight visibility. Every method sees identical trajectories and identical
held-out cells. A dagger marks a matched-budget representative of an architecture
family rather than a reimplementation of the cited paper, as explained in
\cref{sec:experiments}. Adding one geometric channel to a single U-Net is worth more
than any architectural substitution in the panel.}
\label{tab:main}
\end{table}

\subsection{One channel outranks every backbone}

Geometry conditioning outranks every architectural substitution in the panel, and two
comparisons in \cref{tab:main} establish it. The first is that geometry conditioning takes the cascade from $9.99$ to
$8.64$~dB, moving it from behind RadioUNet at $9.93$~dB to ahead of it with
non-overlapping confidence intervals. The second is more surprising. A single
network with the occlusion channel reaches $9.46$~dB and also beats RadioUNet,
while the cascade without the channel does not. The margin there is narrower and
the intervals touch, so we rest it on the paired test rather than on separation of
intervals, but the ordering is clear and the conclusion does not depend on that one
comparison. The architecture is not what separates these methods.

Classical interpolation sits far behind every learned method here, between $25.6$ and
$28.1$~dB, though that gap reflects the trajectory sampling regime rather than a
general ranking and reverses on real measurements. Trajectory-structured sampling
leaves large regions with no nearby observation to interpolate from, which is the
condition classical estimators are least suited to and the one \cref{sec:real}
removes.

\begin{figure}[t]
\centering
\includegraphics[width=\columnwidth]{fig/plateau.pdf}
\caption{\textbf{Architecture moves line of sight and leaves non line of sight
alone.} Open markers are line-of-sight error, filled markers non-line-of-sight
error, one row per method. The nine methods without the occlusion channel spread
over $20.1\%$ in line of sight but only $9.6\%$ in non line of sight, where they fall
inside the shaded band. The two occlusion-conditioned models (red) are the only ones
below it on both axes.}
\label{fig:plateau}
\end{figure}

\subsection{Every deep family stops at the same wall}

Separating the error by visibility shows that architecture moves line-of-sight
accuracy substantially while leaving non-line-of-sight accuracy, which dominates the
total, almost unchanged.
\Cref{fig:plateau} plots line-of-sight against non-line-of-sight error for every
learned method in the panel. Line-of-sight error separates the nine methods
without the occlusion channel over a range of $5.18$ to $9.98$~dB, a coefficient of
variation of $20.1\%$, which is what one expects when comparing convolutional, state-space,
transformer, adversarial and diffusion backbones. Non-line-of-sight error does not
follow. The same nine methods fall between $11.12$ and $15.52$~dB at a coefficient
of variation of $9.6\%$, less than half the variation seen in line of sight, and
dropping PMNet as the one clear outlier tightens the band to $11.12$ to $12.84$~dB
at a coefficient of variation of $5.1\%$. These methods span two orders of
magnitude of design effort.

A plateau that survives this much architectural variation indicates a missing input
rather than insufficient capacity, and supplying that input moves both error modes
below every method in the panel for a few hundred parameters. It is the signature of a quantity no member of the
panel can form from the inputs it is given. Adding \cref{eq:occ} cuts line-of-sight
error to $3.85$~dB and non-line-of-sight error to $9.79$~dB, and the single-network
variant reaches $3.95$ and $10.72$~dB. The channel costs $432$ parameters in the single
network and $576$ in the cascade, increases of $0.006\%$ and $0.009\%$ over
backbones of $7.2$ and $6.4$~million parameters.

The input ablations make the same point in reverse, since withholding a single
geometric input costs more than the entire spread of the architecture panel. Removing
the transmitter heatmap costs $4.76$~dB and removing the building raster costs
$7.32$~dB. The entire learned panel without the occlusion channel, spanning six
architecture families, spreads over $4.18$~dB from best to worst. Withholding one
geometric input therefore moves the result further than every architectural choice
in the comparison put together. What the model is told about the scene dominates
how the model is built.

\begin{figure}[t]
\centering
\includegraphics[width=\columnwidth]{fig/interaction.pdf}
\caption{\textbf{Geometry and depth are substitutes.} Left, the $2\times2$ design.
Right, the gains measured separately sum to $3.55$~dB while the joint configuration
delivers $2.86$~dB. The $0.69$~dB shortfall is the interaction, and its sign means
the cascade has less to add once the geometry is supplied.}
\label{fig:interaction}
\end{figure}

\subsection{Geometry substitutes for depth}
\label{sec:interaction}

The occlusion channel and the cascade do not address different deficiencies, because
their separate gains fail to accumulate by a margin that is itself significant. Were
the two independent, \cref{fig:interaction} would show their benefits adding. Measured
separately against the plain single network, the channel is worth $2.04$~dB and the
cascade is worth $1.51$~dB, for a combined expectation of $3.55$~dB. Applied
together they deliver $2.86$~dB. The interaction term is $+0.69$~dB with a paired
bootstrap interval of $[+0.50, +0.91]$ and $p = 4.7 \times 10^{-10}$ against zero, so
a fifth of the expected joint benefit does not materialise. Each of the four
single-factor comparisons is significant on its own at $p < 10^{-35}$.

Partial redundancy is the only reading consistent with that shortfall, and it implies
that a fixed budget is better spent computing the right input than adding a second
network. A second network in
series has the depth and the intermediate estimate to recover some of the shadow
structure that \cref{eq:occ} supplies directly, so once the structure is given the
cascade has less left to contribute. Stated the other way, geometry buys more where
there is less architecture to spend, which is why the single network gains more
from the channel than the cascade does. A practitioner with a fixed budget is
better served by computing the right input than by stacking a second network.

\subsection{Removing the transmitter does not work}
\label{sec:txfree_result}

The measurement-anchored field of \cref{eq:occmeas} recovers almost none of what the
transmitter-anchored field delivers, even though a paired test can detect the
difference it makes. Substituting $\tilde{O}$ for $O$ in
the cascade gives $9.89$~dB against $8.64$~dB for the transmitter-anchored field,
giving back $1.25$~dB of the $1.35$~dB the channel was worth. What remains is
$0.10$~dB, which a paired test does detect at $p = 2.3 \times 10^{-3}$ over the $320$
scenes but which amounts to $7\%$ of the benefit the transmitter-anchored field
delivers. The substitution is measurable and useless.

The minimum over anchors is what destroys the field, and the reason matters, because
it places the benefit in the transmitter's direction rather than in the building
layout. Taking a minimum over anchors in \cref{eq:occmeas} lets the result be
dominated by whichever measurement happens to sit closest to an
unobstructed path, so the field collapses toward zero almost everywhere. Its mean
value over the scene is $0.084$ against $0.298$ for the transmitter-anchored field,
and the radial shadow structure that carries the signal is gone. The failure is
informative rather than incidental. It says the benefit comes specifically from
knowing where the energy originated, not from knowing the building layout, and any
route to a transmitter-free geometry channel has to preserve directionality rather
than marginalise over it.
